\section{Antécédent électronique et réalisation dimensionnelle}
\label{sec:electron}

La construction électronique reçoit l'état arithmétique avec mémoire de trajectoire produit par APP--TRIT--TPK et ses coordonnées numériques. À ce stade du développement, \(\pi,\varphi,e\) et \(\alpha\) désignent les coordonnées déjà obtenues de la même généalogie, selon l'ordre constructif exposé précédemment. Ce ne sont pas des paramètres tirés d'une table physique. La section électronique est une réalisation ultérieure de cette structure discrète conjointe du continu : elle conserve la cellule d'origine, les deux feuillets arithmétiques, l'orientation, la retenue et la mémoire de retour avant de produire un scalaire.

Il convient de distinguer trois objets dès le départ. Le premier est la section
centrale et son histoire de transport ; le deuxième, la coordonnée
adimensionnelle qui en est obtenue ; le troisième, sa réalisation sur une
droite d'énergie ou de masse. L'égalité entre deux coordonnées numériques
n'identifie pas ces trois niveaux. En particulier, le nombre d'Euler \(e\),
la coordonnée électronique \(\varepsilon_e\) et une énergie de repos
\(E_e\) porteront des notations différentes.

L'objectif de cette section est d'établir la composition
\begin{equation}
 \varepsilon_e^{\beta}
 =\frac{\sqrt3}{4}
  \left[\exp\!\left(\frac{\varphi}{\pi^2}\right)-22\alpha^3\right]
  (1+15\Delta_4)
  R_{\mathrm{act}}^{\,80-54\alpha+6\Delta_4},
 \label{eq:el-formula-completa}
\end{equation}
en expliquant séparément l'origine du préfacteur, les deux entiers de
l'étape basale, la normalisation de \(\Delta_4\), le rapport d'action et le
triplet de retour. La coordinatisation dimensionnelle n'est ouverte qu'après
cette composition. Le résultat ne consiste pas à faire acquérir à un nombre
sans dimension des unités par sa seule ressemblance avec une grandeur
tabulée.

\paragraph{Identité de la section et composition de ses observables.}
\label{el-ampl-articulacion}
L'identité structurelle de l'électron est introduite par la section
centrale et ses opérations de transport. L'inversion de la coordinatisation APP
sélectionne la cellule $(5,5)$ par sa propriété de point fixe. Les deux feuillets
arithmétiques permettent d'étudier les déformations qui conservent la somme et
le reste du produit. Le résultat distingue les deux orientations
$(8,2)$ et $(2,8)$ et conserve la différence d'une unité de quotient.
Tel est le premier niveau d'information : une cellule, une involution, une
fibre conservée et un défaut entier minimal. Aucune évaluation de masse
n'est nécessaire pour l'énoncer ou le démontrer.

Le passage aux projections arithmétiques introduit un deuxième niveau. Les
fibres de somme et de produit déterminent des espérances conditionnelles sur le
même espace fini. Leur commutateur mesure la dépendance à l'égard de
l'ordre de ces deux projections. La norme $\sqrt3/4$ résulte du spectre
de leur appariement ; le bloc qui atteint cette norme conserve en outre
les relations quadratiques démontrées ci-dessous. Le
préfacteur scalaire et l'algèbre locale ont donc une provenance commune, mais
contiennent des informations différentes. La norme retient la grandeur du
commutateur, tandis que les matrices conservent son orientation et sa loi de
composition.

Le transfert surfacique--volumique incorpore les coordonnées régionales
à un troisième niveau. Le calendrier des modes détermine la matrice
d'incidence et son action quadratique. La marque de clôture produit le rapport
$\pi/\varphi^2$ ; l'échange des deux arguments fournit l'
orientation $\varphi/\pi^2$ utilisée dans le caractère exponentiel
électronique. Les deux expressions sont reliées par une involution
explicite. Cette relation explique pourquoi une coordonnée régionale peut
apparaître dans l'équation électronique avec une puissance ou dans une position
différente de celle qu'elle occupait dans la première lecture : l'opération appliquée
fait partie de la donnée mathématique conservée.

Le terme cubique d'alpha appartient au complémentaire régional d'un retour
marqué. Le dénombrement de ses positions produit $27-5=22$, tandis que les
trois feuillets de chacune des cinq régions produisent $3\cdot5=15$.
Les démonstrations conservent le domaine de ces dénombrements et l'
orientation qui détermine leurs signes. La formule basale compose
donc une norme d'incompatibilité, un caractère de
transfert, un déficit de positions et une mémoire torsionnelle.
L'écriture sur une seule ligne exprime la composition de ces
opérations ; son explication exige de les reconstruire dans cet ordre.

La transition de la coordonnée basale à la coordonnée complète conserve
la même section centrale. Le facteur de retour utilise les registres
de l'histoire, et les deux opérateurs de lecture déclarés restituent le triplet
$(80,54,6)$ sur cette trajectoire. La multiplication par
$R_{\mathrm{act}}^{80-54\alpha+6\Delta_4}$ modifie l'évaluation,
mais maintient identifiés son origine et ses antécédents. La valeur
basale constitue une étape intermédiaire de la composition ; la valeur
complète inclut également cette contribution du retour. La comparaison
physique doit utiliser l'étape spécifiée par la formule évaluée.

L'échelle d'action intervient ensuite avec une autre fonction. Le rapport de
deux sections d'action est sans dimension et annule leur échelle commune.
La réalisation absolue utilise, en revanche, une base d'action, un
opérateur de lecture temporel et l'application vers l'énergie ou la masse. Cette distinction permet
de conserver simultanément la construction de l'échelle et l'identité
de la section électronique. Le quantum chronologique fournit une
référence d'échelle ; la section centrale détermine le facteur structurel
exprimé par rapport à elle. Les équations dimensionnelles de ce
chapitre maintiennent cette composition explicite.

La pertinence de l'électron dans l'argument général provient de cette
convergence d'opérations. La même arithmétique qui a produit les
régions détermine maintenant une cellule centrale, ses déformations et ses
projections. La constante de structure fine, construite auparavant, agit
dans cette composition comme coordonnée de clôture. L'information de
retour détermine l'évaluation complète, et la réalisation dimensionnelle
permet la comparaison ultérieure. Expliquer l'électron à partir de la structure
consiste à conserver cette succession d'objets, d'applications et de valeurs à chaque
étape de la démonstration.


\subsection{Centre, orientation et vacance sur une fibre conservée}

Le support cellulaire considéré est \(\{1,\ldots,9\}^2\), avec les deux
évaluations APP, \(r+c\) et \(rc\). L'inversion centrale est
\begin{equation}
 \rho(r,c)=(10-r,10-c).
 \label{eq:el-inversion-central}
\end{equation}
Le centre n'est pas choisi en fonction d'une masse : il se caractérise par une propriété du
support antérieure à sa lecture physique.

\begin{proposition}[Unicité du centre]
L'inversion \(\rho\) a un unique point fixe, \((5,5)\).
\end{proposition}
\begin{proof}
Les équations \(r=10-r\) et \(c=10-c\) équivalent à
\(2r=2c=10\). Leur unique solution dans la coordinatisation nonadique est \((5,5)\).
\end{proof}

La propriété précédente identifie une cellule, et non un historique profond
unique. Une route ajoute le choix initial d'orientation et ses
prolongements compatibles. Dans la coordinatisation utilisée ici, la section
orientée centrale commence en \((5,5)\) avec les deux orientations
positives ; elle sera notée \(\Gamma_e\). Le retour de sa projection
cellulaire n'efface ni les orientations transportées ni le registre de
retenue de la section avec sa mémoire de trajectoire.

La vacance est maintenant considérée, après la clôture de \(\alpha\), comme
une opération locale de cette section. Son arithmétique peut être isolée dans un
lemme fini ; cet isolement n'en fait pas un substitut du
générateur antérieur de constantes.

\begin{definition}[Fibre centrale et défaut de quotient]
La fibre additive centrale est
\[
 \mathcal F_{10}
 =\{(5+t,5-t):t\in\mathbb Z,\ |t|\leq4\}.
\]
Dans la coordinatisation \(n=\rho_9(n)+9q_9(n)\), avec
\(\rho_9(n)\in\{1,\ldots,9\}\), on appelle défaut de quotient par rapport
au centre
\[
 v(t)=q_9(25)-q_9(25-t^2),
\]
lorsque les deux produits ont le même résidu nonadique.
\end{definition}

\begin{lemma}[Vacance centrale minimale]
Les seuls points de \(\mathcal F_{10}\) qui conservent le résidu
multiplicatif du centre sont
\[
 (5,5),\qquad (8,2),\qquad (2,8).
\]
Aux deux points non centraux, le défaut de quotient est exactement un.
\end{lemma}
\begin{proof}
La somme est constamment égale à dix, tandis que
\[
 (5+t)(5-t)=25-t^2.
\]
La conservation du résidu équivaut à \(t^2\equiv0\pmod9\), c'est-à-dire à \(3\mid t\). Dans l'intervalle entier \([-4,4]\), les solutions sont \(t=-3,0,3\). Pour les deux solutions non nulles,
\[
 25=7+9\cdot2,\qquad16=7+9\cdot1.
\]
Le résidu sept est conservé et une unité de quotient est perdue. Il n'y a
pas d'autre déplacement admissible non central dans ce domaine, de sorte que le
défaut positif est minimal. L'échange \(t\mapsto-t\) permute les
deux sorties et conserve le défaut.
\end{proof}

La vacance n'est pas ici un zéro numérique : c'est une différence entre deux
états qui ont la même lecture résiduelle et des quotients distincts. Si l'on
ne conservait que le résidu, cette différence disparaîtrait. L'
orientation distingue en outre les deux points de sortie. Ni le défaut de
quotient ni cette orientation ne s'identifient automatiquement à une charge
électrique ou à une masse ; leur réalisation physique requiert les lecteurs
correspondants.

\subsection{Norme du commutateur des projections arithmétiques}

On obtient le préfacteur électronique sans utiliser \(\alpha\), la torsion ou une unité dimensionnelle. Le domaine est la troisième réalisation d'APP, \(\mathcal A_3=\{1,\ldots,9\}^3\), munie de la mesure uniforme. Nous définissons
\[
 \Sigma(x)=\rho_9(x_1+x_2+x_3),\qquad
 \Pi(x)=\rho_9(x_1x_2x_3).
\]
Les espérances conditionnelles \(P_\Sigma\) et \(P_\Pi\) sont les
projections orthogonales sur les fonctions constantes sur chaque fibre
additive et multiplicative, respectivement. Ainsi, pour une fonction \(f\),
\[
 (P_\Sigma f)(x)
 =\frac{1}{|\Sigma^{-1}(\Sigma(x))|}
   \sum_{y:\,\Sigma(y)=\Sigma(x)}f(y),
\]
et de même pour \(P_\Pi\). Les deux projections agissent sur le même espace ; elles ne représentent pas deux supports arithmétiques indépendants.

\begin{lemma}[Deux projections et angles principaux]
Si \(s\in[0,1]\) est une valeur singulière de l'appariement entre des bases
orthonormales des images de deux projections \(P,Q\), son bloc non
trivial contribue à \(\|[P,Q]\|\) par
\(s\sqrt{1-s^2}\).
\end{lemma}
\begin{proof}
Dans le plan principal correspondant, on peut écrire
\[
 P=\begin{pmatrix}1&0\\0&0\end{pmatrix},\qquad
 Q=\begin{pmatrix}
 s^2&s\sqrt{1-s^2}\\
 s\sqrt{1-s^2}&1-s^2
 \end{pmatrix}.
\]
Leur commutateur est
\[
 [P,Q]=s\sqrt{1-s^2}
 \begin{pmatrix}0&1\\-1&0\end{pmatrix}.
\]
La matrice finale est orthogonale, donc sa norme est un. Les blocs communs ou orthogonaux ont un commutateur nul. La décomposition orthogonale en plans principaux donne l'affirmation et réduit la norme totale au maximum de ces contributions.
\end{proof}

\begin{proposition}[Norme exacte d'incompatibilité]
Sur \(\ell^2(\mathcal A_3)\),
\begin{equation}
 \|[P_\Sigma,P_\Pi]\|=\frac{\sqrt3}{4}.
 \label{eq:el-prefactor}
\end{equation}
\end{proposition}
\begin{proof}
Soit \(N_{ab}=|\Sigma^{-1}(a)\cap\Pi^{-1}(b)|\). L'énumération entière
des trois symboles fournit
\[
 N=9\begin{pmatrix}
 0&1&1&0&1&1&0&1&4\\
 1&0&1&1&0&1&1&0&4\\
 1&0&2&0&1&2&0&0&3\\
 0&1&1&0&1&1&0&1&4\\
 1&0&1&1&0&1&1&0&4\\
 0&0&2&1&0&2&0&1&3\\
 0&1&1&0&1&1&0&1&4\\
 1&0&1&1&0&1&1&0&4\\
 0&1&2&0&0&2&1&0&3
 \end{pmatrix}.
\]
Chaque ligne a pour somme \(81\) ; les sommes des colonnes sont
\[
 (m_1,\ldots,m_9)=(36,36,108,36,36,108,36,36,297).
\]
Les indicatrices normalisées des fibres ont pour matrice
d'appariement \(C_{ab}=N_{ab}/\sqrt{81m_b}\). Par conséquent,
\[
 M=CC^{\mathsf T},\qquad
 M_{ac}=\sum_{b=1}^9\frac{N_{ab}N_{cb}}{81m_b}
\]
est une matrice rationnelle dont le spectre se compose des carrés des
valeurs singulières requises.

La multiplication exacte montre que \((1,\ldots,1)\),
\((1,-1,0)^3\) et \((1,1,-2)^3\) sont des vecteurs propres de valeurs propres
\(1\), \(1/4\) et \(7/132\), respectivement. Le sous-espace des
vecteurs à support dans \(\{3,6,9\}\) dont la somme est nulle a pour dimension
deux et pour valeur propre \(1/18\). Les vecteurs de somme nulle à support dans
\(\{1,4,7\}\) ou dans \(\{2,5,8\}\) forment un noyau de dimension
quatre. Ces sous-espaces sont indépendants et la somme de leurs dimensions vaut
neuf. Le spectre complet est donc démontré
\[
 \operatorname{spec}(M)=
 \left\{1,\frac14,\frac7{132},\frac1{18},\frac1{18},0,0,0,0\right\}.
\]
Le maximum de \(\lambda(1-\lambda)\) sur cet ensemble est \(3/16\), atteint en \(\lambda=1/4\). Le lemme précédent donne \(\sqrt{3/16}=\sqrt3/4\).
\end{proof}

Le mode \((1,-1,0)^3\) est précisément la distribution des valeurs du
sélecteur tritique de phase dans la coordinatisation nonadique. La norme scalaire retient
une mesure d'incompatibilité, mais ne conserve pas à elle seule
l'orientation de ce mode et ne reconstruit pas les deux projecteurs. Cette mémoire
demeure dans l'état de route qui accompagne la lecture scalaire.

\begin{proposition}[Algèbre du bloc électronique local]
Dans le plan principal qui atteint \eqref{eq:el-prefactor}, soient
\[
 P=\begin{pmatrix}1&0\\0&0\end{pmatrix},\qquad
 Q=\begin{pmatrix}1/4&\sqrt3/4\\\sqrt3/4&3/4\end{pmatrix},\qquad
 S=2P-I,\qquad J=\frac4{\sqrt3}[P,Q].
\]
Alors \(S^2=I\), \(J^2=-I\), \(SJ=-JS\), et l'algèbre réelle engendrée par \(S,J\) est \(M_2(\mathbb R)\), réalisation de \(\mathrm{Cl}_{1,1}\). En outre,
\[
 n_\pm=\frac{S\pm J}{2},\qquad
 n_\pm^2=0,\qquad n_+n_-+n_-n_+=I.
\]
\end{proposition}
\begin{proof}
On a \(S=\operatorname{diag}(1,-1)\) et
\(J=\bigl(\begin{smallmatrix}0&1\\-1&0\end{smallmatrix}\bigr)\).
La multiplication vérifie les trois relations. Les matrices \(I,S,J,SJ\) sont linéairement indépendantes et forment une base de \(M_2(\mathbb R)\), ce qui prouve l'affirmation sur l'algèbre. Enfin, \((S\pm J)^2=S^2+J^2\pm(SJ+JS)=0\), et la somme des produits croisés est \((S^2-J^2)/2=I\).
\end{proof}

Dans ce bloc coexistent une direction elliptique, une direction hyperbolique et deux
directions nulles paraboliques. Le résultat décrit une représentation
matricielle locale de l'arithmétique des deux feuillets. Il n'identifie pas toute APP
à \(M_2(\mathbb R)\), ni à un carré latin quantique ; de telles
affirmations exigeraient d'autres domaines et d'autres relations. L'exponentielle
n'appartient pas non plus exclusivement au régime parabolique : elle peut
s'appliquer à n'importe lequel des générateurs de ce bloc.

\subsection{Transfert entre feuillets et coefficients de l'étape basale}

Le calendrier tritique primitif visite les modes
\((\Sigma,\Pi,\Pi)\). En comptant les couples consécutifs, y compris la
fermeture du bloc, on obtient une transition \(\Sigma\to\Pi\), une
\(\Pi\to\Pi\) et une \(\Pi\to\Sigma\). Leur incidence primitive est
\[
 F_{\mathrm{av}}=\begin{pmatrix}0&1\\1&1\end{pmatrix}.
\]
Cette réalisation matricielle s'obtient après le calendrier. Elle n'est pas employée pour remplacer le lecteur coinductif qui a déjà généré \(\varphi\). Son polynôme caractéristique est \(X^2-X-1\), et la reconnaissance des deux racines comme \(\varphi\) et \(-\varphi^{-1}\) décrit l'action de cette sortie dans le transfert modal. L'action quadratique a pour valeurs propres \(\varphi^2,-1,\varphi^{-2}\) ; la dernière est le mode stable positif. La marque régionale \(\pi\), appliquée une seule fois, donne \(\pi/\varphi^2\).

\begin{lemma}[Échange des deux orientations]
Soient \(\mathcal R(x,y)=x/y^2\) et \(J(x,y)=(y,x)\), pour \(x,y>0\).
Alors
\begin{equation}
 J^2=I,\qquad
 \mathcal R(\pi,\varphi)=\frac{\pi}{\varphi^2},\qquad
 (\mathcal R\circ J)(\pi,\varphi)=\frac{\varphi}{\pi^2}
 =\frac1{\varphi^3(\pi/\varphi^2)^2}.
 \label{eq:el-bisagra}
\end{equation}
\end{lemma}
\begin{proof}
L'échange appliqué deux fois est l'identité. De plus,
\(\varphi^3(\pi/\varphi^2)^2=\pi^2/\varphi\) ; prendre son inverse prouve
la dernière égalité.
\end{proof}

Le lecteur électronique utilise l'orientation \(\varphi/\pi^2\),
suivie du caractère exponentiel. L'identité précédente explique
l'échange, mais ne remplace pas la règle qui choisit ce caractère ni ne produit
à elle seule les entiers de la correction. Ceux-ci ont une autre provenance.

Le retour central compatible avec la phase a pour longueur minimale vingt-sept. Avec le registre de phase et de feuillet, son espace fini est
\[
 \mathcal K_e=\mathbb Z/9\mathbb Z\times\mathbb F_3.
\]
Les cinq régions de \(\pi\) sont des identités généalogiques distinctes,
et non cinq copies interchangeables d'un développement visible. Dans la jauge
régionale fixée par la construction, leur injection dans le retour possède
les emplacements
\begin{equation}
 (4,0),\quad(5,0),\quad(6,0),\quad(6,1),\quad(6,2).
 \label{eq:el-cinco-lugares}
\end{equation}
L'orientation transversale, le retour muté et les trois régions positives restent séparés par ce registre. Les changements admissibles de jauge transportent les étiquettes conjointement ; ils ne sélectionnent pas à nouveau les positions à partir d'une valeur électronique.

\begin{proposition}[Déficit et mémoire retenue]
Une fois fixés le retour central et l'injection régionale précédente, le déficit
orienté de frontière et la mémoire torsionnelle sont
\begin{equation}
 n_e=-|\mathcal K_e\setminus\iota(\mathscr R_\pi)|=-22,
 \qquad
 m_e^{\mathrm{tor}}=|\mathscr R_\pi\times\mathbb F_3|=15.
 \label{eq:el-selector-enteros}
\end{equation}
\end{proposition}
\begin{proof}
Les cinq emplacements de \eqref{eq:el-cinco-lugares} sont distincts. Par
conséquent, \(|\mathcal K_e|=9\cdot3=27\) et le complémentaire de l'image
a pour cardinal \(27-5=22\). La convention d'orientation enregistre comme
négatif le déficit retiré de la frontière. La mémoire retenue possède
trois feuillets par région, de sorte que son cardinal positif est
\(5\cdot3=15\).
\end{proof}

La même paire possède le contrôle entier de la coordinatisation APP de la collection :
\begin{equation}
 396=18+72+108+198,
 \qquad \frac{396}{18}=22,
 \qquad \frac{270}{18}=15=3\cdot5=\binom62.
 \label{eq:el-censo}
\end{equation}
Le poids \(18\) appartient au noyau central \(2\times2\), et \(270\) au canal global déclaré. Ces égalités contrôlent la normalisation du décompte ; l'argument de retour fixe en outre les signes. Le quotient des anneaux alternés, \(396/(18+108)=22/7\), est rationnel et ne s'identifie pas à \(\pi\). De même, \(\binom62\) ne signifie pas \(6/2\) : il compte les paires non ordonnées d'un ensemble de six éléments.

L'évaluation du déficit dans la troisième réalisation symétrique utilise
\(\alpha^3\). Cette puissance est le produit de trois copies du scalaire
de clôture déjà obtenu ; elle n'affirme pas de règle universelle faisant
correspondre toute dimension à la même puissance de \(\alpha\).
Le domaine et le lecteur fixés sont essentiels à la composition.

\subsection{La torsion globale et la coordonnée basale}

\begin{definition}[Normalisation torsionnelle électronique]
On distingue le défaut réduit \(d_4\) et le lecteur complet
\(\Delta_4\) :
\begin{equation}
 d_4=\frac{\pi-e\log\pi}{270},\qquad
 \Delta_4=d_4\left(1+\frac{\pi}{729}\right).
 \label{eq:el-delta}
\end{equation}
Le dénominateur \(270\) correspond au canal du recensement, tandis que le facteur \(1+\pi/729\) conserve le couplage avec la cellule cubique de \(729\) positions. La lecture électronique utilise \(\Delta_4\) aussi bien dans son étape basale que dans l'exposant de retour.
\end{definition}

Cette définition agit sur \(\pi\) et \(e\) préalablement générés. Le
logarithme est une opération ultérieure dans leur réalisation réelle, et non une
entrée rétrospective de valeurs conventionnelles. La torsion est globale :
elle ne se détermine pas indépendamment pour chaque section au moyen d'une masse
observée.

\begin{lemma}[Les deux normalisations ne sont pas interchangeables]
On a
\begin{equation}
 \Delta_4-d_4=\frac{\pi}{729}d_4.
 \label{eq:el-delta-diferencia}
\end{equation}
En particulier, \(d_4>0\) et \(\Delta_4>d_4\) pour les coordonnées
considérées.
\end{lemma}
\begin{proof}
L'identité est distributive. Pour la positivité, la fonction
\(f(x)=x-e\log x\), sur \(x>0\), a pour dérivée \(1-e/x\) et
pour dérivée seconde \(e/x^2>0\). Son unique minimum est en \(x=e\), où
elle vaut zéro. Comme \(\pi\ne e\), on a \(f(\pi)>0\). Les deux facteurs
de \eqref{eq:el-delta} sont positifs et \(1+\pi/729>1\).
\end{proof}

\begin{remark}[Conservation de la normalisation]
Supprimer le second facteur de \eqref{eq:el-delta} change le lecteur ; ce n'est pas une simplification algébrique. Si l'on introduit \(d_4\) comme coordonnée auxiliaire, le même objet doit s'écrire \(\Delta_4=(1+\pi/729)d_4\) à chacune de ses occurrences. Maintenir les entiers \(15\) et \(6\) et remplacer seulement \(\Delta_4\) par \(d_4\) produit une autre formule.
\end{remark}

\begin{definition}[Coordonnée électronique basale]
La composition de l'incompatibilité, du transfert modal, du déficit et de la mémoire retenue
est
\begin{equation}
 \varepsilon_e^{(0)}
 =\frac{\sqrt3}{4}
  \left[\exp\!\left(\frac{\varphi}{\pi^2}\right)-22\alpha^3\right]
  (1+15\Delta_4).
 \label{eq:el-basal}
\end{equation}
\end{definition}

L'équation réunit trois opérations de natures différentes. La norme
\(\sqrt3/4\) quantifie l'incompatibilité des projections des
feuillets ; son origine est un angle principal des fibres d'APP.
L'exponentielle agit sur l'orientation \(\varphi/\pi^2\) du
lecteur de transfert. La soustraction cubique enregistre les \(22\) positions
complémentaires du retour marqué, et le facteur torsionnel conserve les
\(15\) composantes correspondant à trois feuillets par région. La
composition conserve ces domaines même si son évaluation finale est scalaire.

La présence de \(\pi\), \(\varphi\) et \(e\) possède ainsi des significations
opératoires. Le nombre d'Euler intervient par l'exponentielle et par
la normalisation logarithmique de \(\Delta_4\) ; \(\varphi\) et \(\pi\)
interviennent dans le transfert orienté. La constante \(\alpha\)
agit après sa propre construction au moyen du terme cubique
de clôture. Le retour chronologique doit encore opérer sur cette coordonnée
basale. La séparation des étapes permet d'expliquer ce qui change lorsqu'on modifie
une orientation, une normalisation ou un registre, tout en conservant la structure
centrale qui identifie la section électronique.

Les parenthèses de cette équation font partie de son contenu. Le terme \(-22\alpha^3\) est soustrait \emph{en dehors} de l'exponentielle ; la mémoire \(1+15\Delta_4\) multiplie toute l'expression entre crochets. Introduire le déficit dans l'exposant ou lui appliquer une torsion différente altérerait la composition.

\begin{proposition}[Positividad basal]
Pour \(0<\alpha<10^{-2}\), \(\varphi>0\) et la normalisation \eqref{eq:el-delta}, on a \(\varepsilon_e^{(0)}>0\).
\end{proposition}
\begin{proof}
L'exponentielle est supérieure à un, et
\(22\alpha^3<22\cdot10^{-6}<1\). Le crochet est donc positif.
On a également \(1+15\Delta_4>1\) et \(\sqrt3/4>0\).
\end{proof}

L'évaluation ultérieure de \eqref{eq:el-basal} commence par
\[
 \varepsilon_e^{(0)}=0.5109989224880660725\ldots.
\]
Ce nombre est encore une coordonnée sans dimension et ne constitue pas la dernière étape : il reste à appliquer la mémoire de retour de la route elle-même.

\subsection{Deux lecteurs du retour électronique}

La mémoire électronique est représentée dans un mot tritique de longueur \(108\) et dans un mot signé de douze phases. Pour que le calcul de l'exposant soit reproductible, nous explicitons la projection chronologique utilisée. Soient initialement \((r_0,c_0)=(5,5)\) et \((h_0,v_0)=(1,1)\). À chaque pas, le sélecteur répète \((+1,-1,0)\). Avec une orientation fixe au sein de chaque bloc de vingt-sept pas, la règle visible est
\[
 \begin{array}{c|c|c}
 \text{mode}&\text{entier lu}&\text{transport cellulaire}\\\hline
 +1&r+c&c\mapsto1+(c-1+h)\bmod9\\
 -1&rc&r\mapsto1+(r-1+v)\bmod9\\
 0&9&\text{sans déplacement}
 \end{array}
\]
Le symbole émis est \(\rho_9(\text{entier lu})\bmod3\).
Après les étapes \(27,54,81\), on inverse simultanément
\(h\) et \(v\). Cette règle calcule la projection observable ; la
mise à jour de l'état conserve en outre le résidu, le quotient,
la retenue, le feuillet, la frontière et la mémoire. Le registre de
\(108\) symboles n'est pas identifié à l'état complet.

Dans chaque bloc complet de trois étapes, on a avancé une fois dans chaque
coordonnée ; neuf de ces blocs sont nécessaires pour revenir au centre dans
la même phase. La longueur minimale de ce retour enregistré est, par
conséquent, \(27\). Si l'on oublie la phase, la cellule est déjà atteinte après
l'étape \(26\), avant le seuil neutre qui complète le bloc ; ce n'est pas
le retour typé que compte \(\mathcal K_e\). La lecture explicite
des deux premiers blocs de retour donne
\begin{equation}
 a=100000220,\qquad b=120200000,\qquad
 \gamma_e=(a^3b^3)^2.
 \label{eq:el-palabra}
\end{equation}
Ici, les puissances désignent la concaténation de mots. En particulier,
\(\gamma_{e,t+54}=\gamma_{e,t}\). Dans le registre de neuf étapes par
phase, les orientations centrales donnent
\begin{equation}
 \tau_e=(1,1,1,-1,-1,-1,1,1,1,-1,-1,-1).
 \label{eq:el-tau}
\end{equation}
La périodicité de ces représentations n'équivaut pas à la réinitialisation de l'histoire qui les produit.

\begin{definition}[Lecteur des discordances cycliques]
Pour un mot \(\gamma\in\mathbb F_3^{108}\), on définit
\[
 D_k(\gamma)=\#\{t\in\mathbb Z/108\mathbb Z:
                   \gamma_t\ne\gamma_{t+k}\}.
\]
Dans le domaine des mots de période divisant \(54\), le lecteur entier est
\begin{equation}
 K_D(\gamma)=\left(D_{27}+D_{36},D_1,\frac{D_4-D_3}{2}\right).
 \label{eq:el-kd}
\end{equation}
\end{definition}

\begin{lemma}[Intégralité et évaluation centrale]
Les \(D_k\) du domaine précédent sont pairs. Pour
\eqref{eq:el-palabra},
\[
 \begin{array}{c|rrrrr}
 k&1&3&4&27&36\\\hline
 D_k&54&56&68&48&32
 \end{array}
\]
et \(K_D(\gamma_e)=(80,54,6)\).
\end{lemma}
\begin{proof}
L'application \(t\mapsto t+54\) apparie sans points fixes les indices de discordance et conserve leur condition ; chaque cardinal est pair. Pour le mot explicite \(a^3b^3\) de longueur \(54\), le décompte des discordances avec chaque décalage est, dans le même ordre, \(27,28,34,24,16\). Doubler le mot double ces cinq cardinaux. Enfin,
\[
 (D_{27}+D_{36},D_1,(D_4-D_3)/2)
 =(48+32,54,(68-56)/2)=(80,54,6).
\]
\end{proof}

Les déplacements \(27\) et \(36\) sont les relèvements
chronologiques des canaux de \(90\) et \(120\) degrés dans une coordinatisation
où neuf pas correspondent à trente degrés. La lecture angulaire
explique cette représentation, mais les entiers sont calculés sur le
mot, sans introduire une masse ni une valeur réelle recherchée.

\begin{definition}[Lecteur de fenêtres dodécaphasiques]
Pour \(\tau\in\{-1,0,1\}^{12}\), avec des indices cycliques, soient
\begin{align*}
 C_k(\tau)&=\sum_{r=0}^{11}\tau_r\tau_{r+k},\\
 A_\ell(\tau)&=\sum_{r=0}^{11}
           \left(\sum_{j=0}^{\ell-1}\tau_{r+j}\right)^2,\\
 P_6(\tau)&=\sum_{r=0}^{5}\tau_r\tau_{r+6}.
\end{align*}
On définit
\begin{equation}
 \Omega(\tau)=4A_3(\tau)-3A_4(\tau),\qquad
 K_\Omega(\tau)=(\Omega(\tau),54,P_6(\tau)).
 \label{eq:el-komega}
\end{equation}
Le \(54\) de ce lecteur enregistre le demi-tour chronologique déclaré.
Il ne se confond pas avec le saut local de valeur quatre d'APP.
\end{definition}

\begin{lemma}[Contraste primitif et forme de corrélation]
Parmi les contrastes linéaires entiers \(uA_3+vA_4\) qui s'annulent
chaque fois que \(A_3=3a\) et \(A_4=4a\), le couple primitif, d'
orientation \(90-120\), est \((u,v)=(4,-3)\). De plus,
\begin{equation}
 \Omega=-2C_1-4C_2-6C_3.
 \label{eq:el-omega-corr}
\end{equation}
\end{lemma}
\begin{proof}
L'annulation pour tout \(a\) exige \(3u+4v=0\). Par coprimalité, les solutions sont \((u,v)=n(4,-3)\), et l'orientation fixe le signe du générateur primitif. En développant les carrés,
\[
 A_\ell=\ell C_0+2\sum_{k=1}^{\ell-1}(\ell-k)C_k.
\]
Par conséquent,
\(A_3=3C_0+4C_1+2C_2\) et
\(A_4=4C_0+6C_1+4C_2+2C_3\). Leur combinaison \(4A_3-3A_4\)
élimine \(C_0\) et produit \eqref{eq:el-omega-corr}.
\end{proof}

L'unicité démontrée appartient à cette classe de contrastes linéaires
entiers ; elle n'exclut pas à elle seule toute fonctionnelle non linéaire des fenêtres.
La spécification de la classe fonctionnelle évite de transformer une identité
de deux canaux en une affirmation de sélection universelle.

\begin{theorem}[Triplet de retour de la section électronique]
Les deux lecteurs appliqués aux registres de \(\Gamma_e\) coïncident et donnent
\begin{equation}
 K_D(\gamma_e)=K_\Omega(\tau_e)=(80,54,6).
 \label{eq:el-triple}
\end{equation}
Son évaluation aux coordonnées déjà générées est
\begin{equation}
 \kappa_e=80-54\alpha+6\Delta_4.
 \label{eq:el-kappa}
\end{equation}
\end{theorem}
\begin{proof}
Le premier lecteur a été calculé plus haut. Pour \(\tau_e\), le produit
antipodal vaut toujours un, d'où \(P_6=6\). Le dénombrement cyclique donne
\(C_1=4\), \(C_2=-4\), \(C_3=-12\) ; donc
\[
 \Omega=-2\cdot4-4(-4)-6(-12)=80.
\]
De manière équivalente, \(A_3=44\), \(A_4=32\) et \(4\cdot44-3\cdot32=80\). Les deux triplets sont donc \((80,54,6)\). L'évaluation de la fonctionnelle \((A,B,C)\mapsto A-\alpha B+\Delta_4C\) prouve la dernière formule.
\end{proof}

Cette coïncidence correspond à la section centrale. Elle n'identifie pas les deux lecteurs dans tous les domaines de routes : l'un compare des symboles tritiques et l'autre compare des fenêtres d'un registre signé. De même, \(80\) ne s'obtient pas à partir de la période d'un mot attribué au nombre d'Euler. Son origine dans cette composition est le triplet de la route électronique.

\subsection{Composition électronique complète et normalisation}

\begin{definition}[Retour multiplicatif]
L'action de la mémoire de retour sur la coordonnée basale est
\begin{equation}
 \varepsilon_e^\beta
 =\varepsilon_e^{(0)}\exp\!\left(\kappa_e\log R_{\mathrm{act}}\right).
 \label{eq:el-beta}
\end{equation}
\end{definition}

\begin{theorem}[Factorisation électronique positive]
Avec les lecteurs, l'orientation et la normalisation déclarés,
\eqref{eq:el-beta} coïncide avec \eqref{eq:el-formula-completa}. Tous
ses facteurs sont adimensionnels et la coordonnée résultante est positive.
\end{theorem}
\begin{proof}
La norme d'incompatibilité est \eqref{eq:el-prefactor} ; le transfert modal,
le déficit et la mémoire retenue composent \eqref{eq:el-basal}. Le
triplet \eqref{eq:el-triple} produit \eqref{eq:el-kappa}, et le rapport des
deux sections d'action est \eqref{eq:el-ratio}. En substituant ces
trois identités dans \eqref{eq:el-beta}, on obtient exactement
\eqref{eq:el-formula-completa}. L'étape basale est positive, tout comme
l'exponentielle d'un nombre réel. Le facteur commun de dimension d'
action s'est déjà simplifié lors de la formation de \(R_{\mathrm{act}}\).
\end{proof}

L'évaluation de cette composition commence par
\begin{equation}
 \varepsilon_e^\beta
 =0.5109989506900008086082571648\ldots.
 \label{eq:el-beta-decimal}
\end{equation}
Les chiffres sont une évaluation ultérieure de la formule, et non des données utilisées
pour sélectionner \(22\), \(15\) ou \((80,54,6)\). La représentation
numérique du terme basal par un préfixe fini introduit uniquement son
erreur de troncature ; elle ne modifie pas la formule exacte qui le définit.

L'incidence d'une normalisation torsionnelle différente peut être localisée sans la comparer à une mesure. Écrivons \(B=\frac{\sqrt3}{4}[\exp(\varphi/\pi^2)-22\alpha^3]\), \(R=R_{\mathrm{act}}\) et
\[
 \mathcal E(d)=B(1+15d)R^{80-54\alpha+6d}.
\]
Pour les deux défauts de \eqref{eq:el-delta},
\begin{equation}
 \frac{\mathcal E(\Delta_4)}{\mathcal E(d_4)}
 =\frac{1+15\Delta_4}{1+15d_4}
   R^{6(\Delta_4-d_4)}>1.
 \label{eq:el-normalizacion-beta}
\end{equation}
La première inégalité torsionnelle et \(R>1\) prouvent qu'il ne s'agit pas de la même coordonnée. Si l'on ne change que le défaut de l'exposant, en conservant le terme basal complet, le rapport des deux sorties est \(R^{6(\Delta_4-d_4)}\), lui aussi différent de un. Une très petite différence reste une différence de formule ; elle ne peut pas être attribuée à une simple conversion d'unités.

\subsection{Réalisation dimensionnelle : énergie, masse et échelle chronologique}

La coordonnée électronique agit finalement sur une droite dimensionnelle.
Soit \(\mathcal L_E\) une droite d'énergie de base positive
\(\mathcal U_E\). La réalisation est
\begin{equation}
 E_e^{(0)}=\varepsilon_e^{(0)}\mathcal U_E,\qquad
 E_e^\beta=\varepsilon_e^\beta\mathcal U_E,
 \qquad m_e^\beta=\frac{E_e^\beta}{c_{\mathrm{int}}^2}.
 \label{eq:el-dimensional}
\end{equation}
La dernière expression appartient à la droite de masse lorsque
\(c_{\mathrm{int}}\) porte la dimension de vitesse dans la réalisation
déclarée. Dans la coordinatisation comparative \(\mathcal U_E=1\,\mathrm{MeV}\),
\eqref{eq:el-beta-decimal} exprime une énergie en MeV ; la masse
correspondante s'exprime en \(\mathrm{MeV}/c^2\). Il est fréquent
d'abréger les deux grandeurs en prenant \(c=1\), mais cette abréviation ne
doit pas masquer l'application dimensionnelle.

L'échelle d'action admet aussi une réalisation action--horloge. Si
\(t_0\) est une section temporelle positive, l'horloge de \(108\) étapes
définit
\begin{equation}
 \omega_0=\frac{2\pi}{108t_0},\qquad
 E_{\mathrm{clk}}=\hbar_{\mathrm{ret}}\omega_0,\qquad
 m_{\mathrm{clk}}=E_{\mathrm{clk}}c_{\mathrm{int}}^{-2}.
 \label{eq:el-cuanto-reloj}
\end{equation}
L'action multipliée par la fréquence a la dimension d'une énergie. Ce quantum
chronologique est une section dimensionnelle commune ; il n'est pas, par définition,
la section électronique. Le mot de route et ses facteurs ne s'
obtiennent pas à partir de l'uniformité de l'horloge.

\begin{proposition}[Annulation de l'échelle commune]
Le facteur \(10^{-34}\mathcal U_S\) de \eqref{eq:el-secciones-accion} ne fait pas partie de \(\kappa_e\). Sur une même droite d'énergie,
\begin{equation}
 \frac{E_e^\beta}{E_e^{(0)}}=R_{\mathrm{act}}^{\kappa_e}.
 \label{eq:el-razon-energia}
\end{equation}
Le résultat est indépendant du choix de la base dimensionnelle
commune.
\end{proposition}
\begin{proof}
Lorsqu'on forme \(\hbar_{\mathrm{pre}}/\hbar_{\mathrm{ret}}\), le facteur
\(10^{-34}\mathcal U_S\) se simplifie exactement. Lorsqu'on divise les deux
expressions d'énergie de \eqref{eq:el-dimensional}, \(\mathcal U_E\) se simplifie
également, et \eqref{eq:el-beta} donne le rapport indiqué.
\end{proof}

Une coordonnée déjà exprimée en MeV n'est donc pas multipliée de nouveau
par \(10^{-34}\). Si l'on choisit \(E_{\mathrm{clk}}\) comme base de
\(\mathcal L_E\), les coordonnées doivent être transportées par le facteur
de changement de base ; conserver inchangé le chiffre de la coordinatisation d'un MeV
reviendrait à identifier deux bases sans le démontrer. C'est précisément ici
que l'échelle d'action entre dans la réalisation absolue et cesse
d'intervenir dans le correcteur relatif.

Il existe également une réalisation linéaire qui rend visible la séparation
entre horloge et généalogie. Dans
\(\ell^2(\mathbb Z/108\mathbb Z;\mathbb R)\), le décalage
cyclique \(S\) a pour espace fixe
\[
 \ker(S-I)=\mathbb R n_0,\qquad
 n_0=\frac1{\sqrt{108}}\sum_{t=0}^{107}\delta_t.
\]
Si \(g_{\Gamma_e}\) est le symbole de la route dans la linéarisation libre
des généalogies et \(\mathcal U_M\) une base de la droite de masse,
on définit l'application de rang un
\begin{equation}
 \mathfrak C_e:
 \mathbb R n_0\otimes\mathbb R g_{\Gamma_e}\longrightarrow\mathcal L_M,
 \qquad
 \mathfrak C_e(n_0\otimes g_{\Gamma_e})
 =\varepsilon_e^\beta\mathcal U_M.
 \label{eq:el-rango-uno}
\end{equation}
L'application réunit la factorisation précédente une fois ses bases déclarées.
Le vecteur uniforme \(n_0\) ne sélectionne pas \(\varepsilon_e^\beta\) :
cette coordonnée procède d'APP, du retour orienté et de ses lecteurs.
On distingue ainsi une réalisation algébrique du mode nul d'une
identification de l'électron au quantum de l'horloge.

La conclusion interne est la factorisation positive d'une section
centrale, avec ses opérations, sa mémoire et ses normalisations explicites.
L'assignation à un électron physique ajoute la représentation de charge, le
spin, le couplage et la coordinatisation correspondante des unités ; elle ne se
déduit pas exclusivement d'une coïncidence décimale. La comparaison
métrologique est présentée après le prolongement exceptionnel et ne modifie
les entrées d'aucune des constructions réalisées ici.

\subsubsection{Dépendance structurelle de la coordonnée électronique}

L'équation électronique exprime une composition hiérarchique. Son premier
facteur, \(\sqrt3/4\), appartient à la structure matricielle centrale et fixe
une normalisation géométrique. L'exposant \(\varphi/\pi^2\) utilise les
coordonnées régionales déjà engendrées dans l'orientation indiquée de l'opérateur
de lecture surfacique--volumique. Le terme \(22\alpha^3\) incorpore le couplage
dodécaphasique à la correction cubique. Le facteur \(1+15\Delta_4\) conserve
la normalisation torsionnelle globale. Enfin, le quotient des sections
d'action agit avec l'exposant déterminé par les opérateurs de lecture du retour.
Chaque opération possède ainsi un antécédent distinct au sein de la même
construction.

L'articulation devient plus précise lorsqu'on considère ses dépendances.
La norme centrale peut être calculée avant l'évaluation des coordonnées
régionales. L'exponentielle et la correction cubique requièrent ces
coordonnées et la clôture de \(\alpha\). La mémoire multiplicative requiert
en outre les sections d'action et la trajectoire de \(108\) positions.
Connaître l'expression basale n'équivaut donc pas à avoir évalué la
coordonnée complète. La succession d'opérations n'ajoute pas une masse
externe au point \((5,5)\) ; elle développe les opérateurs de lecture attribués à sa
section persistante.

La distinction entre échelle et structure fournit également une
interprétation de la comparaison physique. La masse est la réalisation
dimensionnelle de la coordonnée complète, tandis que le centre, l'involution
et la lacune décrivent les relations qui la produisent. Un changement commun
d'unité modifie les coordonnées dimensionnelles, mais laisse inchangés
les rapports sans dimension et les identités de transport démontrées.
C'est en ce sens que l'électron conserve son caractère structurel
lorsqu'il s'exprime par rapport à une échelle chronologique.


\newpage
\subsection{Transport lorentzien et représentation électromagnétique}

La section électronique fournit une structure et une coordonnée
avant son application à une interaction. La réalisation lorentzienne
de l'incidence cubique détermine, avec la convention temporelle déclarée,
l'espace des formes dans lequel cette interaction peut s'exprimer.
Pour la formuler, on fixe une réalisation locale orientée \((\mathcal M,g)\)
de dimension quatre, avec des fibres tangentes identifiées au modèle
\(Q_\square\). On spécifie en outre une connexion \(\nabla\) avec
\(\nabla g=0\), une un-forme potentielle \(A\in\Omega^1(\mathcal M)\),
sa courbure \(F=dA\) et une représentation de charge de la section.
Ces données sont les arguments de la représentation physique qui suit.

La dualité de Hodge satisfait \(\star^2=-I\) sur les deux-formes. La séparation électrique et magnétique correspond à un choix temporel dans la décomposition \(1+3\) ; \(F\) et \(\star F\) conservent leur relation au sein du même espace de deux-formes. La neutralité d'une représentation de charge signifie \(q=0\), tandis que la section peut conserver orientation, phase et transport interne. L'absence de couplage électrique dans cette représentation a un contenu distinct de l'absence d'information de phase.

Dans cette réalisation dimensionnelle, soient \(m_e>0\), \(q\) la charge,
\(\gamma(s)\) une trajectoire paramétrée et \(u=\dot\gamma\)
sa quadrivitesse. La force de Lorentz s'exprime par
\begin{equation}
 m_e\nabla_u u^\mu=qF^\mu{}_{\nu}u^\nu.
 \label{eq:lorentz-electron}
\end{equation}
L'antisymétrie de \(F\) implique
\(F_{\mu\nu}u^\mu u^\nu=0\). Par conséquent,
\[
 \frac{d}{ds}g(u,u)=2g(\nabla_u u,u)=0.
\]
La normalisation de la trajectoire est conservée dès lors que \(\nabla\)
est métrique et que \eqref{eq:lorentz-electron} est satisfaite. La masse utilisée
est la réalisation dimensionnelle de la section électronique précédente ;
le paramètre de couplage est identifié dans la convention électromagnétique
adoptée. La déduction concrète du potentiel et d'une configuration
de champ constitue un problème supplémentaire assorti de conditions de frontière
propres. Le résultat exposé ici détermine le domaine, les données
de la représentation et la conservation qu'implique son équation.

L'inversion des routes du TPK et l'orientation d'une ligne électronique
offrent alors une comparaison précise : chacune possède son involution,
son transport et sa représentation. L'antécédent de Wheeler et Feynman
situe historiquement la question ; l'équivalence de deux réalisations
se décide au moyen de leurs applications et conditions de compatibilité.
Cette disposition permet de maintenir le lien entre l'origine arithmétique de l'électron
et la formulation de son interaction, en explicitant la fonction
de chaque étape.

